% Schematic illustration of Poisson thinning and disjoint increments.
% Positions are deliberately illustrative, not measured or simulated arrivals.
\begin{tikzpicture}[x=0.92cm,y=0.86cm,line cap=round]
 \draw[black!55,thick,-{Stealth[length=2.5mm]}] (0,0)--(9.1,0);
 \node[font=\footnotesize,anchor=north] at (0,-0.09) {$0$};
 \draw[black!45,dashed] (4.5,-0.56)--(4.5,1.12);
 \node[font=\footnotesize,anchor=north] at (4.5,-0.09) {$t$};
 \node[font=\footnotesize,anchor=north] at (8.8,-0.09) {$2t$};
 \node[font=\footnotesize,anchor=south] at (2.25,1.22) {increment on $(0,t]$};
 \node[font=\footnotesize,anchor=south] at (6.65,1.22) {increment on $(t,2t]$};
 \foreach \x in {0.55,2.05,3.25,4.05,5.55,7.45,8.35}{
   \draw[MidnightBlue,thick] (\x,-0.12)--(\x,0.12);
   \fill[MidnightBlue] (\x,0.32) circle (2pt);
 }
 % Red marks are a selection from the same original arrival positions.
 \foreach \x in {0.55,3.25,5.55,8.35}{
   \fill[BrickRed] (\x,0.32) circle (2.7pt);
 }
 \node[font=\footnotesize,anchor=west,text=MidnightBlue] at (0.2,-0.88)
 {charcoal: an original arrival};
 \node[font=\footnotesize,anchor=west,text=BrickRed] at (4.7,-0.88)
 {rust: an arrival kept after thinning};
\end{tikzpicture}
